Notation. $W$ is three-dimensional with $U(3)$ acting by its defining representation (Fact \claimref{standard-model.imported-pair-systems-two-chirality-halves}), so $\Lambda^{k}W$ is the colour representation $1, 3, \bar{3}, 1$ for $k = 0, 1, 2, 3$. A polynomial irreducible $F_{(\lambda_{1}, \lambda_{2})}$ of $U(2)$ has slot number $n = \lambda_{1} + \lambda_{2}$ and weak spin $j = (\lambda_{1} - \lambda_{2})/2$; write $M_{j}(n)$ for the multiplicity of weak spin $j$ in $F_{n}$, so $M_{j}(n) = 0$ for $n < 0$. By Theorem \claimref{standard-model.global-form-quotient-determined} the abelian factor acts on $F_{n} \otimes \Lambda^{k}W$ with $Y = (3n - 2k)/6$, and $V$ is the sum over $(n, k, j)$ of $M_{j}(n)$ copies of $(2j + 1, \Lambda^{k}\C^{3})_{(3n - 2k)/6}$. The complex conjugate of $(2j + 1, R)_{Y}$ is $(2j + 1, \bar{R})_{-Y}$, the $SU(2)$ representations being self-conjugate.

(1) A CPT halving gives $\bar{V} \cong \bar{V}_{+} \oplus \bar{V}_{-} \cong V_{-} \oplus V_{+} = V$; conversely if $V \cong \bar{V}$, every non-real isotype $R$ has the multiplicity of $\bar{R}$, and the real isotypes, which here are colour and hypercharge singlets of $Y = 0$ and so occur at $(n, k) = (0, 0)$ or $(2, 3)$, have even multiplicity when $V$ has the form below, so taking one of each conjugate pair and half of each real isotype is a halving. For the classification compare colour triplets with antitriplets. A triplet occurs only at $k = 1$, where $(2j + 1, 3)_{Y}$ has multiplicity $M_{j}(n)$ with $Y = (3n - 2)/6$; its conjugate $(2j + 1, \bar{3})_{-Y}$ occurs only at $k = 2$, with $-Y = (3n' - 4)/6$, so $n' = 2 - n$. Hence $V \cong \bar{V}$ forces $M_{j}(n) = M_{j}(2 - n)$ for all $j, n$, so $M_{j}(n) = 0$ for $n > 2$. The polynomial irreducibles with $n \le 2$ are $F_{(0,0)}$ ($j = 0$), $F_{(1,0)}$ ($j = \tfrac12$), $F_{(2,0)}$ ($j = 1$) and $F_{(1,1)}$ ($j = 0$); $M_{1}(2) = M_{1}(0) = 0$ excludes $F_{(2,0)}$, and $M_{0}(2) = M_{0}(0)$ gives as many copies of $F_{(1,1)}$ as of $F_{(0,0)}$. Conversely, for $F = a\Lambda^{0} \oplus b\Lambda^{1} \oplus a\Lambda^{2}$ the map $(n, k) \mapsto (2 - n, 3 - k)$ sends each isotype to its conjugate with the same multiplicity ($Y \mapsto (3(2 - n) - 2(3 - k))/6 = -Y$, $\Lambda^{3 - k}\C^{3} \cong \overline{\Lambda^{k}\C^{3}}$, and $M_{j}(n) = M_{j}(2 - n)$ holds), and the real isotype $(1,1)_{0}$ has multiplicity $2a$.

(2) Let $a = 1$. The isotypes come in six conjugate pairs: three from $n \in \{0, 2\}$, namely $(1,3)_{-1/3} \leftrightarrow (1,\bar{3})_{1/3}$, $(1,\bar{3})_{-2/3} \leftrightarrow (1,3)_{2/3}$, $(1,1)_{-1} \leftrightarrow (1,1)_{1}$, each of multiplicity one; two from $n = 1$, namely $(2,1)_{1/2} \leftrightarrow (2,1)_{-1/2}$ and $(2,3)_{1/6} \leftrightarrow (2,\bar{3})_{-1/6}$, each of multiplicity $b$; and the real $(1,1)_{0}$ of multiplicity two. A halving puts $m$ copies of $R$ and $b - m$ of $\bar{R}$ into $V_{+}$ for each pair; write $\epsilon_{1}, \epsilon_{2}, \epsilon_{3} = \pm1$ for the first three (sign $+$ for the first-named member) and $s_{t} = 2m_{t} - b$, $s_{u} = 2m_{u} - b$ for the last two. Each anomaly sum is odd under $R \mapsto \bar{R}$ and vanishes on $(1,1)_{0}$, so with the per-isotype values of Fact \claimref{standard-model.imported-anomaly-coefficients} (gravitational $\sum Y$: $-1, -2, -1, 1, 1$; $[SU(2)]^{2}U(1)$: $0, 0, 0, \tfrac12, \tfrac12$; $[SU(3)]^{2}U(1)$: $-\tfrac13, -\tfrac23, 0, 0, \tfrac13$; $U(1)^{3}$: $-\tfrac19, -\tfrac89, -1, \tfrac14, \tfrac{1}{36}$, for the first-named members in the order above) the conditions are
\[
-\epsilon_{1} - 2\epsilon_{2} - \epsilon_{3} + s_{t} + s_{u} = 0,\quad s_{t} + s_{u} = 0,\quad \epsilon_{1} + 2\epsilon_{2} = s_{u},\quad -\tfrac{\epsilon_{1}}{9} - \tfrac{8\epsilon_{2}}{9} - \epsilon_{3} + \tfrac{s_{t}}{4} + \tfrac{s_{u}}{36} = 0.
\]
The first two give $\epsilon_{3} = -\epsilon_{1} - 2\epsilon_{2}$, which is $\pm1$ only if $\epsilon_{1} = -\epsilon_{2}$, and then $\epsilon_{3} = -\epsilon_{2}$, $s_{u} = \epsilon_{2}$, $s_{t} = -\epsilon_{2}$; the fourth then reads $\epsilon_{2}(4 - 32 + 36 - 9 + 1)/36 = 0$. (As a linear system in five unknowns the four conditions have rank four, with solutions the multiples of $(-1, 1, -1, -1, 1)$.) $s_{u} = 2m_{u} - b = \pm1$ needs $b$ odd, and then $m_{u} = (b \pm 1)/2$, $m_{t} = (b \mp 1)/2$: one generation and $(b - 1)/2$ vectorlike pairs of each doublet type. The $[SU(3)]^{3}$ coefficient $\epsilon_{1} - \epsilon_{2} + 2s_{u}$ vanishes on these solutions and the number of doublets in $V_{+}$, $b + 3b = 4b$, is even.

(3) $F_{1} = a^{\dagger}(\C^{2} \otimes F_{0})$ is $0$ or $\C^{2}$, and $F_{2}$ is a quotient of $\C^{2} \otimes F_{1}$, which is $\mathrm{Sym}^{2}\C^{2} \oplus \Lambda^{2}\C^{2} = F_{(2,0)} \oplus F_{(1,1)}$ (Pieri). If $V$ is CPT-closed, (1) with $\dim F_{0} = 1$ gives $a = 1$, so $F_{2} \supseteq \Lambda^{2}\C^{2} \ne 0$, hence $F_{1} = \C^{2}$, $b = 1$, $F_{2} = \Lambda^{2}\C^{2}$ and $F_{n} = 0$ for $n \ge 3$: $F \cong \Lambda^{\bullet}\C^{2}$. Conversely $\Lambda^{\bullet}\C^{2}$ is generated by the wedge and is CPT-closed by (1). The map $u \otimes u' \mapsto a^{\dagger}(u)a^{\dagger}(u')\lvert 0 \rangle$ is $U(2)$-equivariant from $\C^{2} \otimes \C^{2}$ to $F_{2} = \Lambda^{2}\C^{2}$, so by Schur's lemma it vanishes on $\mathrm{Sym}^{2}\C^{2}$; on $F_{1}$ and $F_{2}$ the product lands in $F_{3} = F_{4} = 0$. So $a^{\dagger}(u)a^{\dagger}(u') + a^{\dagger}(u')a^{\dagger}(u) = 0$ on all of $F$.

(4) For $F = \Lambda^{\bullet}\C^{2}$ ($a = b = 1$), (2) leaves the two solutions $\epsilon_{2} = \pm1$, each fixing $V_{+}$ on every non-real isotype. For $\epsilon_{2} = +1$, $V_{+}$ contains $(1,\bar{3})_{-2/3}$ ($n = 0$, $k = 2$), $(1,\bar{3})_{1/3}$ ($n = 2$, $k = 2$), $(1,1)_{1}$ ($n = 2$, $k = 0$), $(2,3)_{1/6}$ ($n = 1$, $k = 1$) and $(2,1)_{-1/2}$ ($n = 1$, $k = 3$), exactly the components with $n + k$ even other than $n = k = 0$; $\epsilon_{2} = -1$ gives the components with $n + k$ odd other than $(n, k) = (2, 3)$. On the real isotype $(1,1)_{0}$ of multiplicity two, $V_{+}$ is any line, and the anomaly sums do not see the choice. For $\pi = 0$ the grading $(-1)^{k}$ puts $(1,\bar{3})_{-2/3}$ ($n = 0$, $k = 2$) and $(2,\bar{3})_{-1/6}$ ($n = 1$, $k = 2$) into one half, so that half has $s_{u} = -\epsilon_{2}$ in the notation of (2), and its complement likewise; no solution has $s_{u} = -\epsilon_{2}$.

The script enumerates every $F$ with multiplicities at most one among the ten irreducibles with $\lambda_{1} \le 3$, confirming that exactly the three nonzero modules of the form of (1) are CPT-closed, and that adding any of $F_{(2,0)}, F_{(2,1)}, F_{(2,2)}, F_{(3,0)}$ to $a\Lambda^{0} \oplus b\Lambda^{1} \oplus a\Lambda^{2}$ ($a, b \le 3$) breaks closure; it computes with GA.field_kernel the $U(2)$-equivariant maps $\C^{2} \otimes \C^{2} \to \Lambda^{2}\C^{2}$ and $\to \mathrm{Sym}^{2}\C^{2}$ (one each, the first killing symmetric tensors, the second antisymmetric ones), enumerates all generated modules up to slot number four by the Pieri rule and finds $\Lambda^{\bullet}\C^{2}$ the only CPT-closed one, and confirms that truncated bosonic modules are never closed; it counts the anomaly-free halvings for $a = 1$, $b = 0, \ldots, 5$ ($0, 2, 0, 2, 0, 2$) with the cubic and doublet-parity conditions automatic, verifies the coefficient table above from the four sums, the rank-four system and its solution line, the content at $b = 3$, and the identification in (4) of the two solutions with the parity halves, the free neutral line, and the anomalous halves of the $\pi = 0$ grading.
